\documentclass[11pt]{article}
\usepackage[margin=1in]{geometry}
\usepackage{amsmath,amssymb,amsthm}
\newtheorem{theorem}{Theorem}[section]
\newtheorem{lemma}[theorem]{Lemma}
\newtheorem{corollary}[theorem]{Corollary}
\newtheorem{remark}[theorem]{Remark}
\newcommand{\ip}[2]{\langle #1,#2\rangle}
\title{The Quadratic Relaxation Is Sharp Exactly Up to One Radian}
\author{}\date{October 5, 2026 --- paper-proof draft}
\begin{document}\maketitle
\begin{abstract}
The fixed-angle maximizing-cap majorization extends to every $0<\omega<\pi/2$. We extend the strict quadratic equality argument to that whole interval, but prove that the relaxed cap has a strictly larger niche than its signed corner region as soon as $\omega>1$. Consequently the unrestricted optimum is strictly below $1+\omega^2/2$ for $1<\omega<\pi/2$. Together with the exact small-angle theorem, this locates the end of the sharp regime at exactly one radian. It does not identify the maximizing sofa in the remaining interval.
\end{abstract}

\section{The relaxed cap at arbitrary angle}
Let $0<\omega<\pi/2$, $u_t=(\cos t,\sin t)$, $v_t=(-\sin t,\cos t)$,
\[
 c=\sec\omega-\tan\omega,\quad o=(c,1),\quad r=\omega+c-1,
\]
and use the usual fan $F_\omega$. Retain the formulas
\begin{align*}
 p_*(t)&=\omega-t+\ip o{u_t},&k_*(t)&=t+\ip o{v_t},\\
 A(t)&=o+(\omega-t)u_t-v_t,&C(t)&=o+t v_t-u_t.
\end{align*}
Let $K_1=\operatorname{conv}(\{O,o\}\cup A([0,\omega])\cup C([0,\omega]))$. This is Baek's cap $K_{\omega,1}$; the subscript here names the relaxation, not the rotation angle.

\begin{lemma}[The same support formulas remain valid]\label{lem:cap}
$K_1$ is a normalized cap with active supports $p_*,k_*$. These supports are at least one. Its canonical corner
\[
 \gamma(t)=o+(\omega-t-1)u_t+(t-1)v_t
\]
bounds, together with the two radial fan segments, a convex region $D$ satisfying
\[
 \operatorname{int}D\subset N_\omega(K_1),\quad |D|=I(\gamma),\quad
 \mathcal A_1(K_1)=1+\omega^2/2.
\]
\end{lemma}
\begin{proof}
The support calculation in \texttt{relaxation.tex} extends beyond one radian. Its only angle-restricted cross-arc estimate can be replaced by the following. For $s<t\le\omega<\pi/2$,
\[
 s\tan(t-s)\le s\cot s\le1,
\]
where the endpoint $s=0$ is handled directly and $\tan s\ge s$ is used for positive $s$. Thus
$\ip{C(s)-o}{u_t}=s\sin(t-s)-\cos(t-s)\le0$. When $t\le s$ the same inequality is immediate from the signs. The self-arc inequality follows by integrating $A'(t)=(\omega-t)v_t$, and reflection gives the other support interval. Also $p_*'=-1+\cos t-c\sin t\le0$ and $p_*(\omega)=1$, so $p_*\ge1$, and likewise $k_*\ge1$. The same boundary chain and pinned edges therefore describe the cap.

The corner derivative is $\gamma'=-t u_t+(\omega-t)v_t$. Its first coordinate strictly decreases, and its tangent angle increases from $\pi/2$ to $\pi+\omega$, with derivative $1+\omega/(t^2+(\omega-t)^2)>0$. Its ordinate is strictly concave, since
\[
 \gamma_y''=-(t+1)\cos t-(\omega-t+1)\sin t<0.
\]
The endpoint ordinates are $0$ and $r\cos\omega>0$, so the curve lies above the floor. Reflection gives the other fan inequality. Together with the two radial segments it forms the convex region $D$, as in the earlier note.

Both thresholds $p_*(t)-1,k_*(t)-1$ are positive in the interior. Every interior point of $D$ is a strict radial contraction of some $\gamma(t)$ with $0<t<\omega$, and therefore satisfies both strict niche inequalities. This proves the inclusion. Green's formula gives $|D|=I(\gamma)$. The boundary integration proving $|K_1|-I(\gamma)=1+\omega^2/2$ in \texttt{relaxation.tex} is algebraic and uses no restriction $\omega\le1$, proving the last assertion.
\end{proof}

\section{Strict relaxation rigidity on the full open angle interval}
\begin{lemma}[The one-ended Poincare inequality]\label{lem:poincare}
For $f\in H^1([0,\omega])$ with $f(\omega)=0$,
\[
 \int_0^\omega f^2\le\frac{4\omega^2}{\pi^2}\int_0^\omega f'^2.
\]
\end{lemma}
\begin{proof}
Put $\lambda=\pi/(2\omega)$ and $\phi(t)=\cos(\lambda t)$. For continuously differentiable $f$ vanishing at $\omega$, expansion and integration by parts give
\[
 \int_0^\omega(f'^2-\lambda^2f^2)
 =\int_0^\omega\phi^2\bigl((f/\phi)'\bigr)^2\ge0.
\]
The boundary term is zero at $0$ because $\phi'(0)=0$, and at $\omega$ by the first-order zeros of $\phi$ and $f$. The identity may first be taken on truncated intervals and then passed to the endpoint. Approximation in $H^1$, preserving the zero endpoint trace, proves the inequality for the stated class.
\end{proof}

\begin{theorem}[Unique normalized relaxed maximizer]\label{thm:relaxation}
For every $0<\omega<\pi/2$, the unique normalized cap maximizing $\mathcal A_1$ is $K_1$, with value $1+\omega^2/2$. If $f=p-p_*$ and $g=k-k_*$, then
\[
 \mathcal A_1(K_1)-\mathcal A_1(K)
 \ge\frac12\left(1-\frac{4\omega^2}{\pi^2}\right)\int_0^\omega f'^2
       +\frac12\int_0^\omega(g'+f)^2.
\]
\end{theorem}
\begin{proof}
The support-coordinate formula and stationary-point calculation of \texttt{relaxation.tex} apply at every such angle; they use only $f(\omega)=g(0)=0$ and the candidate's endpoint derivative conditions. They give the exact identity
\[
 \mathcal A_1(K_1)-\mathcal A_1(K)
 =\frac12\int_0^\omega\bigl(f'^2+(g'+f)^2-f^2\bigr).
\]
Apply Lemma~\ref{lem:poincare}. Its first resulting coefficient is strictly positive when $\omega<\pi/2$. Equality forces $f'=0$, hence $f=0$ by its endpoint value, then $g'=0$ and $g=0$. The active supports determine the cap. The value follows from Lemma~\ref{lem:cap}.
\end{proof}

\section{An exact negative result beyond one radian}
\begin{theorem}[Extra niche area appears immediately]\label{thm:tails}
For every $1<\omega<\pi/2$,
\[
 |N_\omega(K_1)|>I(\gamma),\qquad
 \mathcal A_\omega(K_1)<\mathcal A_1(K_1)=1+\omega^2/2.
\]
\end{theorem}
\begin{proof}
The intercept of the inner $b$-wall on the horizontal fan ray is
\[
 W(t)=\frac{p_*(t)-1}{\cos t}.
\]
A direct rearrangement gives
\[
 W(t)-r=\frac{(\omega-1)(1-\cos t)+\sin t-t}{\cos t}.
\]
The numerator divided by $t^2$ tends to $(\omega-1)/2>0$ as $t\downarrow0$. Thus $W(t)>r$ for some sufficiently small interior $t$.

Choose $x$ strictly between $r$ and $W(t)$. The point $(x,0)$ satisfies the first niche inequality by $x<W(t)$, and the second because $-x\sin t<k_*(t)-1$, whose right side is positive. A small positive-area rectangle above this point, with all first coordinates greater than $r$, therefore lies in the niche. But $D\subset\{x\le r\}$: the first coordinate of $\gamma$ decreases from $r$, and both radial segments have first coordinates at most $r$. The rectangle misses $D$.

Lemma~\ref{lem:cap} gives $\operatorname{int}D\subset N_\omega(K_1)$ and $|D|=I(\gamma)$. The extra rectangle proves a strictly positive area excess. Subtracting from the cap area yields the second inequality.
\end{proof}

\begin{corollary}[Exact end of the sharp regime]\label{cor:strict}
For the unrestricted fixed-net-angle problem,
\[
 m(\omega)=1+\omega^2/2\quad(0\le\omega\le1),\qquad
 m(\omega)<1+\omega^2/2\quad(1<\omega<\pi/2).
\]
\end{corollary}
\begin{proof}
The equality and uniqueness statement on $[0,1]$ is \texttt{optimality-uniqueness.tex}. For every $0<\omega<\pi/2$, compactness and a positive competitor give a global cap maximizer $K$. The certificate in \texttt{all-angle-bootstrap.tex}, the fixed-angle curvature theorem and the horizontal-monotonicity majorization give
\[
 \mathcal A_\omega(K)\le\mathcal A_1(K)\le1+\omega^2/2.
\]
If equality held at $\omega>1$, Theorem~\ref{thm:relaxation} would force $K=K_1$, contradicting Theorem~\ref{thm:tails}. Thus the cap maximum is strictly below the relaxed value. Every sofa is contained in its fixed-angle monotone envelope, whose area is a cap's area functional. The sofa supremum is therefore bounded by this strict cap maximum. Attainment at the cap level is what makes the strict supremal conclusion valid.
\end{proof}

\begin{remark}[What the negative result does not say]
Non-sharpness of the relaxed value does not by itself prove that the actual cap-minus-niche sofa from $K_1$ is suboptimal; it proves that its area is smaller than the relaxed value. Nor does it identify a new optimal contact regime. An explicit optimizer and its uniqueness for $1<\omega<\pi/2$ require an additional sharp argument. The all-maximizer majorization and strict relaxation kernel are now available, but the positive tail-area correction cannot be discarded.
\end{remark}
\end{document}
